\documentclass[../../../include/open-logic-section]{subfiles}

\begin{document}

\olfileid{sth}{ordinals}{ordtype}
\olsection{ترتيبي عددونه د ترتيب د ډولونو په توګه}

اوس چې تعويض لرو او په $\ZFminus$ کښې کار کوو، بالاخره هغه نتيجه
ثابتولے شو چې مقصد مو وه:

\begin{thm}\ollabel{thmOrdinalRepresentation}
هر ښه ترتيب له يوه يوازيني ترتيبي عدد سره آيزومورف دے.
\end{thm}

\begin{proof}
فرض کړئ $\tuple{A, <}$ ښه ترتيب دے. د
\olref[basic]{ordisoidentity} له مخې، له زيات نه زيات يوه ترتيبي عدد
سره آيزومورف دے. نو د تناقض د راوستلو دپاره فرض کړئ $\tuple{A, <}$ له
\emph{هېڅ} ترتيبي عدد سره آيزومورف نۀ دے. لومړے به «$\tuple{A, <}$ تر
ممکن حده کوچنے کړو». په تفصيل سره: که کومه خاصه پيلنۍ برخه $\tuple{A_a,
<_a}$ له هېڅ ترتيبي عدد سره آيزومورف نۀ وي، نو يو تر ټولو کوچنے $a \in A$
دا خاصيت لري؛ بيا $B = A_a$ او $\mathord{\lessdot} =
\mathord{<_a}$ واخلو. که داسې نۀ وي، نو $B = A$ او $\mathord{\lessdot} =
\mathord{<}$ واخلو.

له تعريفه، د $B$ هره خاصه پيلنۍ برخه له کوم ترتيبي عدد سره آيزومورف
ده، او د پورته بحث له مخې هغه يوازينے دے. نو د تعويض له مخې لاندې
سټ شتون لري او تابعه ده:
\[
	f = \Setabs{\tuple{\beta, b}}{b \in B\text{ او }\ordeq{\beta}{\tuple{B_b, \lessdot_b}}}
\]
د تناقض د بشپړولو دپاره به وښيو چې $f$ يو آيزومورفيزم
$\alpha \to B$ دے، د کوم ترتيبي عدد $\alpha$ دپاره.

ښکاره ده چې $\ran{f}
= B$ دے. او د \olref[iso]{lemordsegments} له مخې $f$ ترتيب ساتي،
يعنې $\gamma \in \beta$ هله او يوازې هله چې $f(\gamma) \lessdot f(\beta)$ وي.
د دې د ښودلو دپاره چې $\dom{f}$ ترتيبي عدد دے، د
\olref[basic]{corordtransitiveord} له مخې دا ښودل کافي دي چې
$\dom{f}$ متعدي دے. نو $\beta \in \dom{f}$ ثابت وټاکئ، يعنې
$\ordeq{\beta}{\tuple{B_b, \lessdot_b}}$ د کوم $b$ دپاره.
که $\gamma \in
\beta$ وي، نو $\gamma \in \dom{f}$ د
\olref[iso]{wellordinitialsegment} له مخې دے؛ په تعميمولو سره، $\beta \subseteq
\dom{f}$ دے.
\end{proof}

دا نتيجه د لاندې تعريف جواز راکوي، چې له \olref[vn]{sec} راهيسې مو
وړاندې کول غوښتل:

\begin{defn}
که $\tuple{A, <}$ ښه ترتيب وي، نو د هغه د ترتيب ډول،
$\ordtype{A, <}$، هغه يوازينے ترتيبي عدد $\alpha$ دے چې
$\ordeq{\tuple{A, < }}{\alpha}$ وي.
\end{defn}

همدارنګه، دا تعريف د دوو ښو اصلونو جواز راکوي:

\begin{cor}\ollabel{ordtypesworklikeyouwant}
که $\tuple{A, <}$ او $\tuple{B, \lessdot}$ ښه ترتيبونه وي:
\begin{align*}
	\ordtype{A, <} = \ordtype{B, \lessdot}&\text{ هله او يوازې هله چې }\ordeq{\tuple{A, <}}{\tuple{B, \lessdot}}\\
	\ordtype{A, <} \in \ordtype{B, \lessdot}&\text{ هله او يوازې هله چې }\ordeq{\tuple{A, <}}{\tuple{B_b, \lessdot_b}}\text{ د يو مناسب }b \in B
\end{align*}
\end{cor}

\begin{proof}
برابري د \olref[basic]{ordisoidentity} له مخې ده. د دويمې ادعا د
ثابتولو دپاره، $\ordtype{A, <} = \alpha$ او $\ordtype{B,
\lessdot} = \beta$ واخلو، او $f \colon \beta \to \tuple {B, \lessdot}$
زموږ آيزومورفيزم وټاکو. بيا:
\begin{align*}
	\alpha \in \beta&\text{ هله او يوازې هله چې }\funrestrictionto{f}{\alpha} \colon \alpha \to B_{f(\alpha)}\text{ آيزومورفيزم دے}\\
	&\text{ هله او يوازې هله چې }\ordeq{\tuple{A, <}}{\tuple{B_{f(\alpha)}, \lessdot_{f(\alpha)}}}\\
	&\text{ هله او يوازې هله چې }\ordeq{\tuple{A, <}}{\tuple{B_b, \lessdot_b}}\text{ د کوم $b \in B$ دپاره}
\end{align*}
دا د \olref[iso]{ordisounique}،
\olref[iso]{wellordinitialsegment} او
\olref[basic]{ordissetofsmallerord} له مخې دي.
\end{proof}

\end{document}
